\section{Tian conditional identification: complete proofs}
\label{sec:annex-p16-tian-conditional-identification}

\begin{annexitem}{annex:p16.tcc-stochastic-bridge}{thm:p16-tcc-stochastic-bridge}{Private response tables realize each stochastic row}
\begin{proof}
For each observed vertex \(i\), let \(\Xi_i\) be the finite
set of its possible observed-parent and latent-input tuples. Take
a private random table \(R_i=(R_{i,\xi})_{\xi\in\Xi_i}\),
whose entries are independent and have masses
\(\mathbb P(R_{i,\xi}=a)=k_i(a\mid\xi)\). Take these
tables independently across vertices and independently of all original
latent inputs. Define the total equation at \(i\) to return
the table entry indexed by its actual parent and input tuple.

For fixed full \(v,u\), the constraints that all observed
equations return \(v\) specify exactly one table entry at
each vertex. Their joint probability over the private tables is
\(\prod_i k_i(v_i\mid v_{\operatorname{pa}(i)},u_i)\).
Acyclic solution uniqueness identifies satisfaction of these equation
constraints with the actual solution having value \(v\).
Summing over the original latent tuple gives the source's mixture
of products. After intervention on \(T\), its table constraints
are removed; all their unused private tables sum to one. The
same argument gives the mixture of truncated products at assignments
consistent with \(t\), and zero otherwise. Thus every intervention
law is preserved. The added inputs are private, so they create
no bidirected edge, and the observed parent arguments are unchanged.
Conversely use the indicators of total equations as the stochastic
rows to obtain the same two formulas directly. This proof uses
the stated independent-input class and makes no conversion claim
for a differently constrained joint latent law.
\end{proof}
\end{annexitem}

\begin{annexitem}{annex:p16.tcc-component-interface}{thm:p16-tcc-component-interface}{The component interface and the exact Q product}
\begin{proof}
On nonempty connected targets and containing sets, the three source
branches are precisely the earlier c-identify branches, after the
renaming of its target and current random set to \(B,H\).
The first branch includes equality \(B=H\), the second therefore
has \(B\subsetneq H\), and in the third the containing
district exists because \(B\) remains connected in the ancestor
graph. Its inclusions are \(B\subseteq H'\subseteq A\subsetneq H\).
The already supplied component invariant, termination and soundness
proofs apply to these exact operations and full Q-kernel arguments.
The later graph-uniform failure conclusion is the previously proved
component completeness result, not an attribution of that proof
to Tian's 2004 open question.

After \(T=t\), the law on \(V\setminus T\) is its Q
kernel. The ancestor set \(D\) is ancestral in this remaining
graph, so the established ancestral summation identity gives
\(P_t(s,c)=\Sigma_F Q[D]\). The unconditional district product
then gives \(Q[D]=\prod_i Q[D_i]\). Every proper observed
parent of a member of \(D\), if not in \(T\), also has
a directed path to \(S\cup C\) avoiding \(T\), and is
therefore in \(D\). This proves the stated inclusive scope bound.
Each \(D_i\) is connected by bidirected paths in the full
graph and lies in one unique maximal district there. Positivity
of \(p\) gives every extracted full-graph Q kernel and every
subsequently reached kernel positive at each assignment. Thus all
ratios and component input tables used in these phases have their
previously proved support and normalization.
\end{proof}
\end{annexitem}

\begin{annexitem}{annex:p16.tcc-quick-criteria}{thm:p16-tcc-quick-criteria}{The singleton-action test and the unchanged conditioning marginal}
\begin{proof}
For the first criterion, consider a hedge for the joint response
\(S\cup C\) under action \(\{X\}\), if one existed.
Its larger forest contains \(X\), while its common roots do
not. Thus \(X\) has a selected directed child in that forest.
The forest is bidirected-connected, so a bidirected path joins
\(X\) to this child there. Every vertex of the forest reaches
a common root by directed edges, and that root reaches the joint
response in the post-action graph. Hence the whole forest is
within \(\operatorname{an}_G(S\cup C)\). The path contradicts
the hypothesis. No hedge exists for the joint query, so its
positive graph-uniform identification follows from the completed
joint theorem. Normalizing its positive conditioning marginal
identifies the conditional query. If \(X\) is outside the
ancestor set, the same conclusion follows directly by ancestral
reduction to the observed marginal.

For the second criterion, nondescendant invariance of acyclic
solutions gives \(C^t=C\) at every latent outcome. Therefore
\(P_t(c)=p(c)>0\). An identified joint determines its conditional
by division. Conversely an identified conditional multiplied by
this observed marginal determines the joint. These implications
hold for every compatible model, and hence for the stated
graph-uniform identification question.
\end{proof}
\end{annexitem}

\begin{annexitem}{annex:p16.tcc-closure}{thm:p16-tcc-closure}{Least fixed-point closure and the two partitions}
\begin{proof}
Initially \(F_0\) contains exactly the marginalized coordinates
in a failed factor's inclusive scope. At every iteration successful
factors are moved only from \(\mathcal I_1\) to \(\mathcal I_0\),
and marginalized vertices only from \(F_1\) to \(F_0\).
The complement definitions therefore preserve both partitions.
Each repetition caused by nonempty \(B_{\mathrm{new}}\)
strictly increases \(F_0\), so there are at most \(|F|\)
such repetitions and one last factor pass. The process terminates.

At its end every factor left in \(\mathcal I_1\) has scope
disjoint from \(F_0\), since the final pass tested every one.
Every marginalized coordinate in a scope of \(\mathcal N\)
was initially included, and every such coordinate in a scope of
\(\mathcal I_0\) was included before the last test found
\(B_{\mathrm{new}}\) empty. Conversely vertices enter \(F_0\)
only for those two reasons. Hence \(F_0=F\cap H'\).

To prove leastness, take any pair satisfying the same seed and
closure rules. It must contain the initial \(F_0\). If it
contains the current \(F_0\), it must include every successful
factor whose scope meets that set, and then every marginalized
vertex in those factors' scopes. Induction through the iteration
shows that it contains the final pair. The final pair itself
satisfies these rules, so it is the unique least one. In particular
different orders of testing or moving eligible factors give the
same closed pair. All factors in \(\mathcal I_0\) remain
individually identified; their relocation records interaction with
the unknown block rather than loss of their individual formulas.
\end{proof}
\end{annexitem}

\begin{annexitem}{annex:p16.tcc-cancellation-soundness}{thm:p16-tcc-cancellation-soundness}{The two finite sums and valid cancellation}
\begin{proof}
If \(\mathcal N\) is empty, the component theorem identifies
every factor in \(P_t(s,c)=\Sigma_F\prod_i Q[D_i]\).
Summation over \(S\) gives its positive conditioning marginal,
so the Phase-2 return is the conditional effect.

Otherwise fix any compatible model and assume Phase 4 succeeds.
Let
\[
 B(v)=\prod_{D_i\in\mathcal N\cup\mathcal I_0}Q[D_i](v),
 \qquad G_1(v)=\prod_{D_i\in\mathcal I_1}Q[D_i](v).
\]
The first block's arguments lie in \(H'\). The closure theorem
and the success test put its arguments among \(F_0,C,T\):
it has none in \(F_1\) or \(S\). The second block has
no argument in \(F_0\), again by closure; its arguments lie
among \(F_1,S,C,T\). Coordinates outside \(D\cup T\)
enter neither block by the scope lemma. Finite distributivity thus
gives
\[
 P_t(s,c)=\sum_{f_0,f_1}B(v)G_1(v)
       =\left(\sum_{f_0}B(v)\right)
          \left(\sum_{f_1}G_1(v)\right)=U(c,t)K(s,c,t).
\]
This is the source's Eq.~(28) with the scopes and coordinate
binders justified. Each actual Q factor is positive at every
assignment, and each finite sum has at least one term, so both
\(U\) and \(K\) are positive. Summing over response values
gives \(P_t(c)=U(c,t)\sum_{s'}K(s',c,t)\). Cancellation
of the positive \(U\) is therefore valid, and the remaining
denominator is positive. Every factor in \(K\) has a known
component expression, so the returned conditional depends only
on \(p\), the graph and assigned values, in every compatible
model. The proof never assumes that the unknown block itself
is determined from observations. Empty sets and products satisfy
the same calculation; in particular an empty response gives one.
\end{proof}
\end{annexitem}

\begin{annexitem}{annex:p16.tcc-complexity}{thm:p16-tcc-complexity}{Termination and a graph-operation bound}
\begin{proof}
There are at most \(n\) full-graph or ancestral districts.
An ancestor search and a bidirected component search each cost
\(O((n+1)^2)\) with an adjacency matrix. In a component call,
the current random set strictly decreases whenever recursion occurs,
so there are at most \(n+1\) calls. At most \(n\) target
districts require such a computation. All their graph decisions
therefore cost \(O((n+1)^4)\). The finite kernel sums and
products in a successful call are defined by positivity, so they
also terminate when numerical tables are supplied.

In Phase 3, at most \(n\) marginalized vertices and \(n\)
factors can be moved. One complete pass uses at most \(n\)
intersections of scopes of size at most \(n\), followed by
unions of such sets, hence \(O((n+1)^2)\) elementary set
operations in an incidence-array representation. At most \(n+1\)
passes suffice. Initial decompositions and the final intersection
test do not exceed these bounds. This proves the stated polynomial
bound for graphical work. Symbolic sums over named coordinate
sets and products can be retained as shared formula nodes rather
than expanded tables, with polynomially many such construction
operations along these calls. The argument makes no bound on
the number of numeric cells generated by expanding those sums,
nor on exact arithmetic with unrestricted value-set cardinalities
or precision. These are distinct from the source's graphical
polynomial-time assertion.
\end{proof}
\end{annexitem}

\begin{annexitem}{annex:p16.tcc-moral-bridge}{thm:p16-tcc-moral-bridge}{Connected district blocks in the moral and incidence graphs}
\begin{proof}
Work after fixing \(T\), retaining the observed ancestor set
\(D\) and the canonical latent-root expansion. Any latent
root with only one retained observed child is a private leaf
input. Adding one formal private root at every observed vertex
does not change observed d-separation: such a root has degree
one and cannot be internal on a path between observed endpoints.
It also cannot change directed ancestry from an observed collider.
The same statement holds after the required observed arrow deletions.
We include these private roots to make every block below nonempty.

The ancestral graph of \(S\cup C\) contains every observed
vertex of \(D\), all its latent parents, and no additional
observed vertex after fixing \(T\). Moralize this graph and
remove \(C\). For each district \(D_i\), group its latent
parents, its retained observed outputs \(D_i\setminus C\),
and its observed parents outside \(C\cup T\) into a block.
Latent roots with more than one retained child belong to one
district, and private roots are assigned to their child's district.
Blocks may overlap at observed parent/output vertices.

Each block is connected in the remaining moral graph. Along a
bidirected spanning tree inside \(D_i\), successive latent
parents are connected through a shared observed child when that
child is not conditioned; when it is in \(C\), they are
connected by the marriage edge created at that child. Its private
root joins these same parents by marriage. An observed parent
outside \(C\cup T\) joins the child directly if retained,
or joins its latent/private parents by marriage if that child
is conditioned. For a singleton district the private root gives
the same connected block, even if its observed output is removed.
Thus every member of the inclusive scope remaining in \(D\setminus C\)
is attached to this block.

Conversely every edge of the remaining moral graph belongs to
at least one such block. An observed parent-child edge belongs
to the child's district block. An edge from a latent parent
belongs to its unique district block. A marriage edge joins two
parents of one observed child, and belongs to that child's block.
There are no arrows between latent roots in the canonical expansion.
Consequently a chain of overlapping blocks is exactly a path in
the factor-scope incidence graph, and a path in either graph
can be converted to a chain of blocks in the other. This proves
the claimed correspondence of response-connected components.

For a conditioning vertex \(c\), the IDC moral-graph criterion
from the exchange-closure proof says that it is eligible exactly
when none of its parents lies in a response-connected component
after removing all conditioning vertices. Every remaining parent of \(c\)
belongs to its district block. If that block is not response-
connected, the criterion holds. If it is response-connected, its
private formal parent is in that component because the block
is connected, so the criterion fails. The private-root addition
does not change the original separation test, as shown above.
Thus eligibility is exactly absence of \(D(c)\) from
\(\mathcal R\), as claimed.
\end{proof}
\end{annexitem}

\begin{annexitem}{annex:p16.tcc-idc-bridge}{thm:p16-tcc-idc-bridge}{Closure reachability, maximal exchanges and the exact terminal ancestors}
\begin{proof}
Consider the incidence graph with only the marginalized variable
vertices \(F\) retained, and start from the failed factor
vertices \(\mathcal N\). Alternating graph reachability first
finds their marginalized neighbors, then the successful factors
touching those vertices, then their marginalized neighbors, and
so on. These are exactly the seed and movement steps of Phase 3.
The least-closure theorem therefore identifies its reachable
variable vertices with \(F_0\) and its additional factor
vertices with \(\mathcal I_0\). Phase 4 fails precisely if
one of these factors has a response vertex in its scope.
Equivalently, a failed factor connects to \(S\) in the full
incidence graph. Any such path can be stopped at its first
response vertex, so all internal variable vertices are in \(F\).
Thus failure is exactly \(\mathcal N\cap\mathcal R\ne\varnothing\).

The moral bridge and IDC's maximal-exchange theorem give
\(E=C\setminus D_{\mathcal R}\) and remaining conditioning
\(W=C\cap D_{\mathcal R}\). We prove that the terminal
joint ancestor set is exactly \(D_{\mathcal R}\).
This set contains \(S\cup W\). If an observed parent \(u\)
of its vertex lies outside \(T\cup E\), then either it is
in \(C\), hence in \(W\), or it is an unconditioned
vertex of \(D\). In the second case it is adjacent in the
incidence graph to the response-connected factor of its child.
Its own factor also contains it in its inclusive scope, so that
factor is response-connected too. Hence \(u\in D_{\mathcal R}\).
This proves ancestral closure in the graph with \(T\cup E\)
removed, giving one inclusion for the terminal ancestor set.

For the reverse inclusion, take \(v\in D_{\mathcal R}\)
outside \(S\cup W\). It lies in \(F\) and has an
original directed path, avoiding \(T\), to \(S\cup C\).
Stop that path at its first such target. All intermediate
vertices are unconditioned. Starting from \(v\), each edge
to its next child puts that child's factor into \(\mathcal R\),
since the current parent is an unconditioned response-connected
variable. Thus all intermediate vertices lie in \(D_{\mathcal R}\).
If the endpoint is in \(C\), its factor is also in
\(\mathcal R\), so the endpoint belongs to \(W\), not
\(E\). The path therefore reaches \(S\cup W\) without
passing through a newly fixed vertex. This proves the other
inclusion. Vertices already in \(S\cup W\) use length-zero paths.

The set \(D_{\mathcal R}\) is a union of whole maximal
districts of \(G[D]\), so its districts are exactly those
in \(\mathcal R\): no edges join different original districts,
and each retained one remains intact and connected. The terminal
joint identification problem therefore asks for exactly their
Q factors. Its full-graph containing districts and component
subroutine calls are the same calls used in Tian's Phase 2.
The already proved completeness of the unconditional modified
Tian procedure, and its equivalence with ID, show that this
terminal joint ID call fails exactly when one of these factors
belongs to \(\mathcal N\). This is the same criterion as
Phase 4, completing the bridge.
\end{proof}
\end{annexitem}

\begin{annexitem}{annex:p16.tcc-completeness}{thm:p16-tcc-completeness}{The corrected graph-uniform completeness conclusion}
\begin{proof}
Success soundness was proved directly by Tian's two-sum cancellation,
without replacing his algorithm by IDC. The bridge now proves
that its failed-factor closure has the same success/failure
decision as IDC's terminal joint problem after the unique maximal
exchange set. IDC is sound and complete on the stated positive
graph-uniform finite class. Therefore Tian's procedure succeeds
exactly on those conditional queries identified in that class.
Both successful expressions equal the same conditional kernel.

If Tian's procedure fails, the bridge supplies the failed terminal
joint problem and the completed IDC conditional-witness theorem
supplies equal-observation positive finite countermodels for the
original conditional query. This gives the required necessity
through an explicit graphical bridge, rather than by assuming
that an unknown factor cannot cancel numerically. It also explains
the later addendum's closure argument with its inclusive parent
scopes retained. The countermodels establish the graph-uniform
claim; they do not classify each fixed numerical law or each
prescribed small alphabet family separately.
\end{proof}
\end{annexitem}

\begin{annexitem}{annex:p16.tcc-original-examples}{thm:p16-tcc-original-examples}{Every source component call and cancellation partition}
\begin{proof}
\paragraph{Figure 2, conditioning only on W.}
The full-graph districts are \(\{B\},\{Z\}\) and
\(H_1=\{A,X,W,Y\}\). In the topological order
\(A,B,X,W,Z,Y\), their identified kernels are
\(Q[\{B\}]=p(b\mid a)\), \(Q[\{Z\}]=p(z\mid x)\)
and exactly the stated \(q_1\). All these expressions use
positive conditional denominators. After fixing \(X\), the
ancestors of \(Y,W\) are \(D=\{Y,W,Z\}\). Its
districts are the three singletons and \(F=\{Z\}\).

The W call begins with \((\{W\},H_1,q_1)\). Its ancestor
set inside \(H_1\) is \(\{X,W\}\), so it forms
\(q_{XW}=\Sigma_{\{A,Y\}}q_1\). This set is one district
because of \(X\leftrightarrow W\), and its extraction keeps
that whole kernel. The next call has directed \(X\to W\)
and bidirected \(X\leftrightarrow W\); its ancestor set
equals the whole current set and properly contains its target.
It returns failure. The W scope is \(\{X,W\}\).

The Y call begins with \((\{Y\},H_1,q_1)\). Its ancestor
set is \(\{X,W,Y\}\), whose districts are \(\{X,W\}\)
and \(\{Y\}\). The ancestor marginal is
\(\Sigma_{\{A\}}q_1\), and extracting the final Y singleton
in order \(X,W,Y\) gives
\[
 q_Y=\frac{\Sigma_{\{A\}}q_1}{\Sigma_{\{A,Y\}}q_1}.
\]
The following singleton call returns this table at its ancestral
base. The true kernel \(Q[\{Y\}]\) uses outside observed
parents only at \(W,Z\); thus this expression is independent
of its other displayed fixed arguments. The Z singleton was
already a full-graph district and its component call returns
\(p(z\mid x)\) immediately.

Consequently \(\mathcal N=\{\{W\}\}\) and
\(\mathcal I=\{\{Y\},\{Z\}\}\). Its initial bad scope
is \(\{X,W\}\), which misses \(F=\{Z\}\), so
\(F_0=\varnothing\), \(F_1=\{Z\}\) and
\(\mathcal I_0=\varnothing\). No factor is moved. The
bad scope misses response Y, and cancellation returns
\(K(y;w,x)=\sum_z p(z\mid x)q_Y(y;w,z)\) divided by
its sum over Y. Since each \(q_Y\) sums to one over Y
and \(p(z\mid x)\) sums to one over Z, this denominator
is one. The uncancelled joint is
\(Q[\{W\}](w;x)K(y;w,x)\), and its W marginal is
\(Q[\{W\}](w;x)\). This proves the source's explicit
cancellation calculation, including the unknown common factor.

\paragraph{Figure 2, conditioning additionally on B.}
The post-action ancestor set is now \(D=\{A,B,W,Z,Y\}\),
since A is a parent of the additional conditioning vertex B.
Its districts are \(\{A,Y\},\{B\},\{W\},\{Z\}\),
and \(F=\{A,Z\}\). The target \(\{A,Y\}\) is
bidirected-connected. Inside its containing full district \(H_1\),
every vertex is an ancestor of A or Y, so its component call
fails immediately. The W call fails as before, while B and Z
are identified. Thus the failed family contains \(\{A,Y\}\)
and \(\{W\}\). Their inclusive scopes include
\(\{A,Y,W,Z\}\) and \(\{X,W\}\), respectively.
Both marginalized vertices A and Z enter \(F_0\). The B
and Z factors meet this set and enter \(\mathcal I_0\).
The final scope already contains response Y through the failed
\(\{A,Y\}\) factor, so Phase 4 fails. The necessity of
that graph decision follows from the complete bridge theorem;
it is not inferred just from the written appearance of an
unknown symbol in one numerical formula.

\paragraph{Figure 4 and its four phases.}
The full-graph districts are \(\{B\},\{Z\}\), and
\(H_1=\{X,A,W,Y\}\). In order \(X,A,B,W,Z,Y\)
their factors are \(p(b\mid a)\), \(p(z\mid x)\),
and the stated \(q_2\). For the query \(P_x(Y\mid A)\),
the post-action ancestor set is all of \(\{A,B,W,Z,Y\}\),
with districts \(\{A\},\{B\},\{Z\},\{W,Y\}\),
and marginalized set \(F=\{B,W,Z\}\).

For \(\{W,Y\}\), the ancestor set inside \(H_1\) is
just \(\{W,Y\}\), since the observed route from X to W
passes through B, outside that current graph. The base branch
therefore returns \(q_{WY}=\Sigma_{\{X,A\}}q_2\).
This is normalized over \(W,Y\) and has outside observed
parents only \(B,Z\). For \(\{A\}\), the ancestor
set in \(H_1\) is \(\{X,A\}\). Its marginal is
extracted as the one district \(\{X,A\}\); the next call
fails because \(X\to A,X\leftrightarrow A\). The B and
Z calls return their known singleton factors. This gives
\(\mathcal N=\{\{A\}\}\) and successful districts
\(\{B\},\{Z\},\{W,Y\}\).

The initial bad scope \(\operatorname{Pa}^{+}(\{A\})=\{X,A\}\)
misses \(F\). Thus \(F_0\) and \(\mathcal I_0\)
remain empty, the bad scope misses Y, and Phase 4 returns
the stated ratio. All sums over X and A belong only to the
subexpression \(q_2\); they do not bind the action X in
\(p(z\mid x)\) or the evidence A in \(p(b\mid a)\).
Its denominator can also be evaluated as one: sum \(q_{WY}\)
over W,Y first, then the normalized B and Z factors. This
completes every component call and partition step of the source examples.
\end{proof}
\end{annexitem}

\begin{annexitem}{annex:p16.tcc-corrected-comparison}{thm:p16-tcc-corrected-comparison}{The corrected scope computation and explicit positive countermodels}
\begin{proof}
For this example the source notation is \(T=\{X\}\),
\(S=\{W\}\), \(C=\{Z\}\). The post-action ancestor
set is \(D=\{W,Y,Z\}\), marginalized set \(F=\{Y\}\),
and its districts are all three singletons. The full graph's
districts are \(\{X,Z\},\{Y\},\{W\}\).
The W and Y factors are known as \(p(w)\) and \(p(y\mid w)\).
The Z call from \(Q[\{X,Z\}]\) fails on the bow
\(X\to Z,X\leftrightarrow Z\). Hence the failed family
is exactly \(\{\{Z\}\}\).

Its inclusive scope is \(\operatorname{Pa}^{+}(\{Z\})=\{Z,X,Y\}\),
so the initial \(F_0\) is \(\{Y\}\). The successful
Y factor has inclusive scope \(\{Y,W\}\); it intersects
\(F_0\) at its own output coordinate Y and must move into
\(\mathcal I_0\). The W factor has scope \(\{W\}\)
and stays in \(\mathcal I_1\). There are no remaining
marginalized vertices to add. The final bad-block scope is
\(\{Z,X,Y,W\}\), which contains the response W. Phase 4
therefore fails. Omitting the output Y from its own factor's
scope would skip exactly this required propagation step. The
original paper's inclusive definition resolves the comparison
without changing the algorithm.

For the concrete models, naturally \(X=U\), \(W=V\),
\(Y=V\oplus E_Y\), and \(Z=Y\oplus E_Z\) in both.
For every observed assignment the required noise values are
\(e_Y=y\oplus w\), \(e_Z=z\oplus y\). The independent
fair \(U,V\) contribute \(1/4\), giving the stated common
law. Its cells are \(9/64\) when both noise values are zero,
\(3/64\) when exactly one is one, and \(1/64\) when both
are one. For each \((x,w)\), the four corresponding noise
assignments have total mass \((9+3+3+1)/64=1/4\); hence
all sixteen positive cells sum to one.

After \(X=x\), the first model has
\(Z=x\oplus U\oplus Y\oplus E_Z\). Its independent
fair \(U\) makes Z fair and independent of W, so
\(P^1_x(W=1\mid Z=z)=1/2\). In the second model,
\(Z=W\oplus E_Y\oplus E_Z\). The effective flip
probability is \(2(1/4)(3/4)=3/8\). With W fair,
Z is fair and Bayes' formula gives response-one probabilities
\(3/8\) at Z=0 and \(5/8\) at Z=1. Both conditioning
masses are \(1/2\), so this is a valid positive conditional
countermodel pair. IDC also makes no exchange because the
path \(Z\leftarrow Y\leftarrow W\) is active after
the proposed outgoing-arrow deletion at Z; its terminal joint
request has the same failed Z factor. Thus both procedures
correctly fail on the comparison graph.
\end{proof}
\end{annexitem}
